% Part: incompleteness
% Chapter: introduction
% Section: definitions

\documentclass[../../../include/open-logic-section]{subfiles}

\begin{document}

\olfileid{inc}{int}{def}

\olsection{定义}

为了实施 Hilbert 将数学形式化并证明该形式化既一致又完备的计划，首要任务将是选择一门语言、一个逻辑框架以及一套公理系统。为了我们的目的，不妨假设数学可以在一阶语言中形式化，也就是说，存在某个由常项符号、函数符号和谓词符号组成的集合，它们与一阶逻辑的联结词和量词一起，使我们能够表达数学命题。多数人同意这样的语言存在：集合论的语言，其中 $\in$ 是唯一的非逻辑符号。如此简单的语言竟具有如此强的表达力，初看起来当然是一个难以令人信服的断言，而且，要确立几乎所有数学都能用这个极其贫乏的词汇来表达这一事实，需要付出大量的工作。为了简单起见，现在让我们把讨论限制在算术，即只涉及自然数~$\Nat$ 的那部分数学。用来表达算术事实的最自然的语言是~$\Lang L_A$。$\Lang L_A$ 包含一个二元谓词符号~$<$、一个常项符号~$\Obj 0$、一个一元函数符号~$\prime$ 以及两个二元函数符号~$+$ 和~$\times$。

\begin{defn}
一个语句集合~$\Gamma$ 是一个\emph{理论}，当且仅当它对语义后承封闭，即 $\Gamma = \Setabs{!A}{\Gamma \Entails
!A}$。
\end{defn}

有两种简单的指定理论的方法。一种是将其指定为某个结构中为真的语句的集合。例如，考虑 $\Lang L_A$ 的这样一个结构：其论域是~$\Nat$，并且所有的非逻辑符号都按预期的方式解释。

\begin{defn}\ollabel{def:standard-model}
\emph{算术的标准模型}是结构~$\Struct N$，定义如下：
\begin{enumerate}
\item $\Domain N = \Nat$
\item $\Assign{\Obj 0}{N} = 0$
\item $\Assign{\Obj \prime}{N}(n) = n + 1$ 对所有 $n \in \Nat$ 成立
\item $\Assign{\Obj +}{N}(n, m) = n + m$ 对所有 $n, m \in \Nat$ 成立
\item $\Assign{\Obj \times}{N}(n, m) = n\cdot m$ 对所有 $n, m \in \Nat$ 成立
\item $\Assign{\Obj <}{N} = \Setabs{\tuple{n, m}}{n \in \Nat, m \in
  \Nat, n < m}$
\end{enumerate}
\end{defn}

注意 `$\times$' 和 `$\cdot$' 的区别：$\times$ 是算术语言中的一个符号。当然，我们选择它是为了提醒我们乘法，但 $\times$ 并非乘法运算，而是一个二元函数符号（严格来说是 $\Obj f^2_1$）。相比之下，$\cdot$~\emph{才是}通常的乘法函数。当你看到像 $n \cdot m$ 这样的表达式时，我们指的是数字 $n$ 和~$m$ 的乘积；而当你看到像 $x \times y$ 这样的表达式时，我们讨论的是算术语言中的一个项。在标准模型中，函数符号~$\times$ 被解释为自然数上的函数 $\cdot$。对于加法，我们用 $+$ 既表示算术语言的函数符号，又表示自然数上的加法函数。这里需要根据上下文来确定其含义。

\begin{defn}
\emph{真算术}的理论是在算术标准模型中被满足的语句的集合，即
\[
\Th{TA} = \Setabs{!A}{\Sat{N}{!A}}.
\]
\end{defn}

$\Th{TA}$ 是一个理论，因为若 $\Th{TA} \Entails !A$，则 $!A$ 在每个满足~$\Th{TA}$ 的结构中都被满足。由于 $\Sat{N}{\Th{TA}}$，我们有~$\Sat{N}{!A}$，因此 $!A \in \Th{TA}$。

指定理论~$\Gamma$ 的另一种方式，是将其指定为某个语句集~$\Gamma_0$ 的所有语义后承组成的集合。在这种情况下，$\Gamma$ 是 $\Gamma_0$ 在语义后承下的“闭包”。只有当 $\Gamma_0$ 被显式指定时（例如，将~$\Gamma_0$ 的元素一一列出），以这种方式指定一个理论才有意义。至少，$\Gamma_0$ 必须是可判定的，也就是说，必须存在一种可计算的检验方法，以判定某个语句是否属于~$\Gamma_0$。我们称~$\Gamma_0$ 中的语句为~$\Gamma$ 的\emph{公理}，而称~$\Gamma$ 是由~$\Gamma_0$ \emph{公理化}的。

\begin{defn}
一个理论~$\Gamma$ 是由~$\Gamma_0$ \emph{公理化}的，当且仅当
\[
\Gamma = \Setabs{!A}{\Gamma_0 \Entails !A}
\]
\end{defn}

\begin{defn}
由下列语句公理化而成的理论 $\Th{Q}$ 被称为“Robinson 的 $\Th{Q}$”，它是一个非常简单的算术理论。
\begin{align*}
& \lforall[x][\lforall[y][(\eq[x'][y'] \lif \eq[x][y])]] \tag{$!Q_1$}\\
& \lforall[x][\eq/[\Obj 0][x']] \tag{$!Q_2$}\\
& \lforall[x][(\eq[x][\Obj 0] \lor \lexists[y][\eq[x][y']])] \tag{$!Q_3$}\\
& \lforall[x][\eq[(x + \Obj 0)][x]] \tag{$!Q_4$}\\
& \lforall[x][\lforall[y][\eq[(x + y')][(x + y)']]] \tag{$!Q_5$}\\
& \lforall[x][\eq[(x \times \Obj 0)][\Obj 0]] \tag{$!Q_6$}\\
& \lforall[x][\lforall[y][\eq[(x \times y')][((x \times y) + x)]]] \tag{$!Q_7$}\\
& \lforall[x][\lforall[y][(x < y \liff \lexists[z][\eq[(z' + x)][y]])]] \tag{$!Q_8$}
\end{align*}
语句集 $\{!Q_1, \dots, !Q_8\}$ 是 $\Th{Q}$ 的公理，因此 $\Th{Q}$ 由这些公理的所有语义后承组成：
\[
\Th{Q} = \Setabs{!A}{\{!Q_1, \dots, !Q_8\} \Entails !A}.
\]
\end{defn}

\begin{defn}
设 $!A(x)$ 是 $\Lang L_A$ 中的一个公式，其自由变元为 $x$ 和 $y_1$, \dots, $y_n$。那么任何形如
\[
\lforall[y_1][\dots\lforall[y_n][((!A(\Obj 0) \land \lforall[x][(!A(x)
\lif !A(x'))]) \lif \lforall[x][!A(x)])]]
\]
的语句都是\emph{归纳模式}的一个实例。

\emph{Peano 算术}~$\Th{PA}$ 是由 $\Th{Q}$ 的公理连同所有归纳模式的实例公理化而成的理论。
\end{defn}

\begin{explain}
归纳模式的每个实例都在~$\Struct{N}$ 中为真。
如果公式~$!A$ 只含一个自由变元~$x$，这一点最易理解。
此时，$!A(x)$ 在~$\Struct{N}$ 中定义了~$\Nat$ 的一个子集~$X_{!A}$。
$X_{!A}$ 是所有使得在 $s(x) = n$ 时 $\Sat{N}{!A(x)}[s]$ 成立的~$n \in \Nat$ 组成的集合。
对应的归纳模式实例是
\[
((!A(\Obj 0) \land \lforall[x][(!A(x) \lif !A(x'))]) \lif
  \lforall[x][!A(x)]).
\]
如果它的前件在~$\Struct{N}$ 中为真，那么 $0 \in X_{!A}$，并且若 $n \in X_{!A}$，则 $n+1 \in X_{!A}$。
既然 $0 \in X_{!A}$，我们就得到 $1 \in X_{!A}$；由 $1 \in X_{!A}$ 又得到 $2 \in X_{!A}$。依此类推。
所以，对每个 $n \in \Nat$，都有 $n \in X_{!A}$。
但这意味着 $\lforall[x][!A(x)]$ 在~$\Struct{N}$ 中被满足。
\end{explain}

$\Th{Q}$ 和 $\Th{PA}$ 都是公理化的理论。
一个大问题是：它们有多强？例如，$\Th{PA}$ 能否证明所有可以用~$\Lang L_A$ 表达的关于~$\Nat$ 的真理？
具体而言，$\Th{PA}$ 的公理能否解决所有可以用~$\Lang L_A$ 表述的问题？

另一种表述方式是问：是否有 $\Th{PA} = \Th{TA}$？
显然，$\Th{TA}$ 确实能证出（即包含）所有关于~$\Nat$ 的真理，并且能解决所有可以用 $\Lang L_A$ 表述的问题，
因为如果 $!A$ 是 $\Lang L_A$ 中的一个语句，那么要么 $\Sat{N}{!A}$，要么 $\Sat{N}{\lnot !A}$，从而要么 $\Th{TA} \Entails !A$，要么 $\Th{TA} \Entails \lnot !A$。
我们把这样的理论称为\emph{完全的}。

\begin{defn}
一个理论~$\Gamma$ 是\emph{完全的}，当且仅当对其语言中的每个语句~$!A$，要么 $\Gamma \Entails !A$，要么 $\Gamma \Entails \lnot !A$。
\end{defn}

\begin{explain}
由完备性定理，$\Gamma \Entails !A$ 当且仅当 $\Gamma \Proves !A$，
所以 $\Gamma$ 是完全的当且仅当对其语言中的每个语句~$!A$，要么 $\Gamma \Proves !A$，要么 $\Gamma \Proves \lnot !A$。
\end{explain}

我们由此产生的另一个问题是：是否存在一个计算过程，可以用来检验一个语句是否属于~$\Th{TA}$、$\Th{PA}$，或者仅仅属于~$\Th{Q}$？
我们可以通过定义一个集合（例如语句集）何为\emph{可判定的}，使这个问题更加精确。

\begin{defn}
一个集合 $X$ 是\emph{可判定的}，当且仅当存在一个计算过程，该过程对输入~$x$，
若 $x \in X$ 则返回~$1$，否则返回~$0$。
\end{defn}

于是我们的问题变成了：$\Th{TA}$（$\Th{PA}$，$\Th{Q}$）是可判定的吗？

所有这些问题的答案都将是：否。这些理论都是不可判定的。
然而，这种现象并非这些特定理论所独有。事实上，\emph{任何}满足某些条件的理论都适用于同样的结论。
$\Th{Q}$ 和 $\Th{PA}$ 满足的条件之一，就是它们由可判定的公理集所公理化。

\begin{defn}
一个理论是\emph{可公理化的}，如果它由一个可判定的公理集所公理化。
\end{defn}

\begin{ex}
任何由有限语句集公理化的理论都是可公理化的，因为任何有限集都是可判定的。
因此，例如 $\Th{Q}$ 就是可公理化的。

像~$\Th{PA}$ 这样通过模式公理化的理论也是可公理化的。
因为要检验 $!B$ 是否是~$\Th{PA}$ 的公理之一（即计算函数 $\Char{X}$，其中若 $!B$ 是~$\Th{PA}$ 的公理则 $\Char{X}(!B) = 1$，否则 $\Char{X}(!B) = 0$），我们可以这样做：
首先，检查 $!B$ 是否是~$\Th{Q}$ 的公理之一。如果是，答案为“是”且 $\Char{X}(!B) = 1$。
如果不是，则检验它是否是归纳模式的一个实例。
这可以系统地进行；下面或许是最容易看到的一种具体做法：
归纳模式的任何一个实例都以若干全称量词开始，然后是一个作为条件句的子公式。
该条件句的后件是 $\lforall[x][!A(x, y_1, \dots, y_n)]$，其中 $x$ 和 $y_1$, \dots, $y_n$ 是~$!A$ 的所有自由变元，并且~$!B$ 开头的量词约束了变元~$y_1$, \dots,~$y_n$。
一旦我们提取出这个~$!A$，并检查了它的自由变元与开头的那些全称量词以及~$\lforall[x]$ 所约束的变元相匹配，
我们就继续检查该条件句的前件是否与 \[
!A(\Obj 0, y_1, \dots, y_n) \land \lforall[x][(!A(x, y_1, \dots, y_n)
\lif !A(x', y_1, \dots, y_n))]
\] 相匹配。
同样，如果匹配，$!B$ 就是归纳模式的一个实例；如果不匹配，则不是。
\end{ex}

在回答这个问题——以及哪些理论是完全的或可判定的这一更一般的问题——时，考虑以下定义也会有所帮助。
回想一下，集合 $X$ 是\emph{可数的}，当且仅当它为空集或存在一个满射函数~$f \colon \Nat \to X$。这样的函数称为 $X$ 的一个枚举。

\begin{defn}
一个集合 $X$ 被称为\emph{可计算枚举的}（简称 c.e.\@），当且仅当它是空集或它拥有一个可计算的枚举。
\end{defn}

除了可公理化之外，不完全性定理所适用的理论还需满足另一个条件，即它们必须足够强，足以证明关于可计算函数和可判定关系的基本事实。所谓“基本事实”，是指那些表达可计算函数在每个自变元处取什么值的语句。而“足够强”是指，所讨论的理论将这些语句作为其定理。例如，考虑一个典型的可计算函数：加法。$+$ 在自变元 $2$ 和 $3$ 处的值是~$5$，即 $2+3 = 5$。算术语言中表达 $+$ 在自变元 $2$ 和 $3$ 处值为~$5$ 的语句是：$(\num{2} + \num{3}) = \num{5}$。例如，$\Th{Q}$ 就能证明这个语句。更一般地，我们希望对于每个可计算函数 $f(x_1, x_2)$，都存在~$\Lang{L_A}$ 中的一个公式~$!A_f(x_1, x_2, y)$，使得只要 $f(n_1, n_2) =
m$，就有 $\Th{Q}
\Proves !A_f(\num{n_1}, \num{n_2}, \num{m})$。这样，$\Th{Q}$ 就证明了~$f$ 在自变元 $n_1$、$n_2$ 处的值是~$m$。事实上，我们要求它证明得还要多一点，即当自变元为 $n_1$、$n_2$ 时，没有其他数是~$f$ 的值。对于可判定关系也有类似的要求。下面的两个定义将这一点精确化。

\begin{defn}
公式~$!A(x_1, \dots, x_k, y)$ 在~$\Gamma$ 中\emph{表示}函数 $f\colon \Nat^k \to \Nat$，当且仅当若 $f(n_1,
\dots, n_k) = m$，则
\begin{enumerate}
\item $\Gamma \Proves !A(\num{n_1}, \dots, \num{n_k}, \num{m})$，且
\item $\Gamma \Proves \lforall[y](!A(\num{n_1}, \dots, \num{n_k},
y) \lif y = \num{m})$。
\end{enumerate}
\end{defn}

\begin{defn}
公式~$!A(x_1, \dots, x_k)$ \emph{表示}关系 $R \subseteq \Nat^k$，当且仅当：
\begin{enumerate}
\item 若 $R(n_1, \dots, n_k)$，则 $\Gamma \Proves !A(\num{n_1},
\dots, \num{n_k})$，且
\item 若并非 $R(n_1, \dots, n_k)$，则 $\Gamma \Proves \lnot
!A(\num{n_1}, \dots, \num{n_k})$。
\end{enumerate}
\end{defn}

要使不完全性定理适用，一个理论必须是“足够强”的，即它能表示所有可计算函数和所有可判定关系。$\Th{Q}$~及其扩张满足这个条件，但要证明这一点需要我们花些时间——这是关于$\Th{Q}$能证明哪些事实的一个非平凡结论，并且由于$\Th{Q}$仅有少数几条公理，我们必须从中推出所有这些事实，因此证明起来并不容易。然而，$\Th{Q}$是一个非常弱的理论。所以，尽管证明$\Th{Q}$表示所有可计算函数很困难，但大多数有趣的理论都比~$\Th{Q}$更强，也就是说，它们能证明比$\Th{Q}$更多的结论。而如果$\Th{Q}$证明了某个结论，任何更强的理论也能证明它；既然$\Th{Q}$能表示所有可计算函数，那么每个更强的理论也能。这意味着许多有趣的理论都满足不完全性定理的这个条件。因此，我们的努力不会白费，因为这表明不完全性定理适用于广泛的理论。当然，任何旨在形式化“全部数学”的理论都必须证明$\Th{Q}$所能证明的一切，因为它至少应该能够涵盖初等计算的结果。所以，任何作为“全部数学”理论候选者的理论，都将是不完全性定理适用的对象。